\section{Required Future Research and Experiments}
\label{sec:required-future-research-r015}

This chapter records the scientific work that the present release does not claim to have completed. It is part of the claim boundary of Core~1.3.3: where a statement needs more proof, a cleaner certificate, a prospective replay design, or a stronger comparator study, the main text must point here rather than pretend that the gap has already been closed. The rows below are not excuses for missing work. They are the current research frontier after the source-level review gate has demoted unsupported public claims.

\begin{longtable}{@{}L{0.08\textwidth}L{0.19\textwidth}L{0.28\textwidth}L{0.31\textwidth}@{}}
\caption{Required future research and experiments for claims that are bounded, demoted, or not yet promoted in Core~1.3.3.}\\
\toprule
\textbf{ID} & \textbf{Title} & \textbf{Purpose and hypothesis} & \textbf{Method and possible outcomes} \\
\midrule
\endfirsthead
\toprule
\textbf{ID} & \textbf{Title} & \textbf{Purpose and hypothesis} & \textbf{Method and possible outcomes} \\
\midrule
\endhead
FR-001 & Prospective K-Level Prediction Battery & Design a prospective test set for K-level transition hypotheses without using retrospective examples as proof. Goal: Determine which K-level obligations survive independent future cases. Hypothesis: If the K-level hierarchy captures real structural transitions, independently selected cases should preserve the predicted support/demotion pattern under predeclared criteria. Current blocker: A prospective public dataset and pre-registered case battery are not yet available inside the current release resources. & Pre-register observable families, lock source snapshots, run finite/replay checks, compare residuals against negative controls, then classify support, demotion, or falsifier. Possible outcomes: Support would strengthen domain projection; mixed results would refine K-level criteria; failure would demote or restructure the affected K-level claims. \\
FR-002 & Clean Lean Certificate Binding & Repair or rebuild the Lean/source-manifest binding so mechanized support can be promoted only when source hashes and theorem names match. Goal: Separate formalization inventory from machine-checked theorem support without ambiguity. Hypothesis: A clean certificate will either promote a smaller mechanized subset or expose exact theorem obligations that remain prose/finite only. Current blocker: The current certificate records binding failures and cannot be honestly promoted as clean support. & Rebuild the Lean project from the pinned source tree, bind theorem names to proof sheets, record toolchain, source hashes, and failure rows. Possible outcomes: Clean build promotes selected machine-checked rows; mismatch keeps Lean as inventory and creates formal work orders. \\
FR-003 & Prospective Replay And Comparator Study & Convert retrospective bounded replay QA into a prospective or externally held-out replay design. Goal: Test whether replay/comparator rows predict unseen targets rather than reconstructing known rows. Hypothesis: If OC projection constraints add value, locked formulas should outperform specified comparator baselines on held-out rows. Current blocker: No independently locked prospective replay corpus is packaged in Core 1.3.3. & Lock formulas, data snapshots, comparator baselines, negative controls, residual thresholds, and failure rules before evaluation. Possible outcomes: Positive results support domain projection; null or adverse results demote empirical language and refine formulas. \\
FR-004 & All-Domain Scope Boundary Audit & Map where OC can speak as model grammar, where it has theorem support, and where it has empirical support. Goal: Prevent architectural ambition from becoming unsupported all-domain assertion. Hypothesis: The release is strongest as a claim-governed model core until enough proof and replay rows exist per domain. Current blocker: The current corpus does not yet contain complete proof/evidence lanes for every domain projection. & For each domain projection, require a source packet, theorem/proof or finite witness, comparator, replay/evidence row, falsifier, and explicit nonclaim. Possible outcomes: A complete row can be promoted narrowly; an incomplete row remains an example or future-research task. \\
\bottomrule
\end{longtable}
